%%
%% Chapter 1: Introduction
%% Independent file, included by ../main.tex
%%

\section{Introduction}
\label{sec:intro}

Cone-beam computed tomography (CBCT) is a cornerstone of image-guided
procedures in radiation therapy and interventional radiology, offering
volumetric imaging with compact gantry geometries. Clinical workflows,
however, place hard limits on the number of projection views that can be
acquired: fewer views reduce radiation dose and shorten scan times, but
dramatically undersample the Fourier domain of the object. Sparse-view CBCT
reconstruction -- recovery of a high-quality volume from a handful of
projections, typically far below the hundreds required by analytical
reconstruction -- is therefore an ill-posed inverse problem whose solution
requires strong prior knowledge of anatomy.

Deep learning has become the dominant tool for this task. Neural
reconstruction methods range from direct 2D$\to$3D back-projection networks
to score-based and diffusion models, and recent work has shown that
explicit scene representations such as 3D Gaussian splatting (3DGS)
\cite{kerbl2023gaussiansplats} provide an outstanding inductive bias for
tomographic volumes: compact, differentiable, and easily supervised by
projection consistency. The DIF-Gaussian framework \cite{lin2024difgaussian}
demonstrated that a splatting-based reconstructor can produce accurate
volumes from as few as six cone-beam projections.

Despite these advances, the vast majority of sparse-view CBCT models are
\emph{deterministic}: they map a fixed set of views to a single volume and
output no measure of confidence. This is unsatisfactory in a clinical
setting, where (i)~reconstruction noise and ambiguity are spatially
inhomogeneous and task-relevant, and (ii)~downstream planning decisions
(e.g., dose painting or target delineation) should ideally consume
calibrated uncertainty, not a point estimate. Probabilistic reconstruction
addresses this by placing a distribution over latent scene content and
marginalizing it at inference.

A principled probabilistic design, however, must reconcile two competing
goals. On one hand, training can exploit \emph{dense} observations (many
views) that are too expensive to acquire at deployment; conditioning a
variational posterior on these dense views concentrates information and
tightens the evidence lower bound. On the other hand, the deployment-time
input is exclusively the \emph{sparse} set, so the model must not silently
depend on information -- neither dense views nor ground-truth volume -- that
is unavailable at inference. Any mismatch between the training-time and
deployment-time information flow is a latent bug, and, as we show, a
surprisingly common one: it is the direct cause of \emph{posterior
collapse}, in which the learned variational posterior degenerates to the
uninformative prior and the latent variable carries no information.

In this paper we design, build, and rigorously diagnose a probabilistic
extension of the Gaussian-splatting CBCT reconstructor. Our contributions
are:

\begin{enumerate}
\item \textbf{A neural-process-style probabilistic formulation for
splatting-based CBCT reconstruction.} We inject a tractable global latent
variable into the Gaussian-splatting decoder through feature-wise linear
modulation (FiLM), yielding a latent neural process (dense-NP) whose prior
$p(z\mid C)$ is a function of the sparse context views $C$ alone.

\item \textbf{A GT-conditioned CVAE variant (P5).} During training only, a
lightweight 3D CNN encodes the ground-truth volume into the variational
posterior $q(z\mid C,Y)$, giving the model access to dense target
information while the objective simultaneously supervises the
deployment-side prior samples. At inference $q$ is discarded and the model
uses $p(z\mid C)$ exclusively; \emph{no} ground truth or dense views are
required.

\item \textbf{Systematic latent diagnostics and a concrete failure
analysis.} We instrument training with KL tracking, cross-case latent
statistics, prior/posterior reconstruction gaps, and zero- and
shuffle-perturbation interventions. This isolates posterior collapse as a
persistent failure mode of a seemingly reasonable design choice --
conditioning the decoder on dense-view features during training
(\texttt{posterior\_decoder\_views='dense'}) -- and identifies a
deterministic mean-code control as the strongest reference.

\item \textbf{State-of-the-art sparse-view results with calibrated
uncertainty.} On the LUNA16 CBCT benchmark with six input views, our best
GT-conditioned CVAE variant reaches $29.24$\,dB PSNR / $0.863$\,SSIM, within
$0.3$\,dB of the deterministic mean-code reference ($29.54$\,dB /
$0.870$\,SSIM), while additionally producing deployment-side variance
estimates for every reconstructed point.
\end{enumerate}

The remainder of this paper is organized as follows. Section~\ref{sec:rel}
reviews related work on sparse-view CT reconstruction, Gaussian
splatting, and latent variable models. Section~\ref{sec:method} presents
the proposed probabilistic formulation and the P5 algorithm in detail.
Section~\ref{sec:exp} describes the experimental protocol, results, and
diagnostics. Section~\ref{sec:conclusion} concludes and outlines future
directions.